\section{The Klein-four symmetry $K_4 = \langle{\sim},{-}\rangle$}

This section records claim nodes and dependency edges.

\begin{proposition}[Conflation and its composite]\label{thm:conflate}\lean{Obsmetrologia.Charge.conflate}\leanok\uses{thm:four-object}
\texttt{Charge.conflate}, written~$-$, is the knowledge-negation $\langle a, b\rangle \mapsto \langle \neg b, \neg a\rangle$: it swaps $\bot_k \leftrightarrow \top_k$ and fixes $t$ and $f$. Together with negation~$\sim$ it generates $K_4$. The companion \texttt{Charge.negConf} is the coordinatewise complement $\langle a, b\rangle \mapsto \langle \neg a, \neg b\rangle$, the third nontrivial element ${\sim}\circ{-}$.
\end{proposition}
\begin{theorem}[The generators commute]\label{thm:negate-conflate-comm}\lean{Obsmetrologia.Charge.negate_conflate_comm}\leanok\uses{thm:conflate}
\texttt{Charge.negate\_conflate\_comm} proves ${\sim}\circ{-} = {-}\circ{\sim}$: negation and conflation commute. \texttt{Charge.negate\_conflate} identifies their common value as \texttt{negConf}, the coordinatewise complement. Commutation is the hinge that makes the four composites close into a group rather than a larger structure.
\end{theorem}
\begin{theorem}[Negation's order-effects]\label{thm:negate-orders}\lean{Obsmetrologia.Charge.negate_preserves_kle}\leanok\uses{thm:conflate, thm:four-bilattice}
\texttt{Charge.negate\_preserves\_kle} proves $\sim$ preserves the knowledge order ($\mathrm{KLE}\,({\sim}a)\,({\sim}b) \leftrightarrow \mathrm{KLE}\,a\,b$), while \texttt{Charge.negate\_reverses\_tle} proves $\sim$ reverses the truth order ($\mathrm{TLE}\,({\sim}a)\,({\sim}b) \leftrightarrow \mathrm{TLE}\,b\,a$). Negation is the symmetry that fixes knowledge and dualizes truth.
\end{theorem}
\begin{theorem}[Conflation's order-effects]\label{thm:conflate-orders}\lean{Obsmetrologia.Charge.conflate_reverses_kle}\leanok\uses{thm:conflate, thm:kle, thm:four-bilattice}
Dually, \texttt{Charge.conflate\_reverses\_kle} proves $-$ reverses the knowledge order and \texttt{Charge.conflate\_preserves\_tle} proves $-$ preserves the truth order. The two generators partition the two orders' duties between them: each fixes one order and reverses the other.
\end{theorem}
\begin{theorem}[The composite reverses both]\label{thm:negconf-reverses}\lean{Obsmetrologia.Charge.negConf_reverses_kle}\leanok\uses{thm:negate-orders, thm:conflate-orders}
\texttt{Charge.negConf\_reverses\_kle} (with its truth-order companion) proves ${\sim}\circ{-}$ reverses \emph{both} orders. The product of two symmetries each reversing exactly one order reverses both --- the $\mathbb{Z}/2 \times \mathbb{Z}/2$ structure read off the order-effects.
\end{theorem}
\begin{theorem}[The Klein-four group]\label{thm:klein-four-symmetry}\lean{Obsmetrologia.Charge.klein_four}\leanok\uses{thm:negate-conflate-comm, thm:negconf-reverses}
\texttt{Charge.klein\_four} bundles the group facts of $\{\mathrm{id}, \sim, -, {\sim}\circ{-}\}$: each of the three nontrivial elements is an involution, and the products $\sim\circ- = {\sim}\circ{-} = -\circ\sim$ close the set. This is $K_4 \cong \mathbb{Z}/2 \times \mathbb{Z}/2$, the natural symmetry group of FOUR equipped with both negation and conflation --- the two orders' duality made a group.
\end{theorem}

This section records claim nodes and dependency edges.

\begin{proposition}[The carrier $K_4$ and its action]\label{thm:k4-group}\lean{Obsmetrologia.Positing.Gauge.K4}\leanok\uses{thm:klein-four-symmetry}
\texttt{Positing.Gauge.K4} is the four-constructor inductive $\{\texttt{id}, \texttt{neg}, \texttt{conf}, \texttt{negConf}\}$ --- the Klein-four group as a first-class carrier. Its action \texttt{Positing.Gauge.K4.act} sends each element to the corresponding map on \texttt{Charge} (\texttt{neg}~$\mapsto\sim$, \texttt{conf}~$\mapsto-$, \texttt{negConf}~$\mapsto{\sim}\circ{-}$), and \texttt{Positing.Gauge.K4.mul} is the explicit $4\times 4$ Cayley table.
\end{proposition}
\begin{theorem}[The Cayley table]\label{thm:cayley}\lean{Obsmetrologia.Positing.Gauge.K4.cayley}\leanok\uses{thm:k4-group}
\texttt{Positing.Gauge.K4.cayley} records the full $4\times 4$ multiplication grid, one \texttt{rfl} per cell. It fixes the group's combinatorics before any law is proven from it.
\end{theorem}
\begin{theorem}[The action is a homomorphism]\label{thm:act-mul}\lean{Obsmetrologia.Positing.Gauge.K4.act_mul}\leanok\uses{thm:cayley, thm:k4-group}
\texttt{Positing.Gauge.K4.act\_mul} proves $\mathrm{act}\,(g \cdot h)\,a = \mathrm{act}\,g\,(\mathrm{act}\,h\,a)$: the Cayley table is precisely composition of the underlying \texttt{Four} maps. This is what makes $K_4$ a group \emph{acting on} FOUR rather than an abstract table, and it is here that the involution and commutation laws of the symmetry are consumed.
\end{theorem}
\begin{theorem}[$K_4$ is an abelian group object]\label{thm:k4-isgroup}\lean{Obsmetrologia.Positing.Gauge.K4.isGroup}\leanok\uses{thm:act-mul}
\texttt{Positing.Gauge.K4.isGroup} bundles associativity (\texttt{mul\_assoc}), the two-sided identity \texttt{id}, the self-inverse law \texttt{mul\_self} ($g \cdot g = \texttt{id}$ for every $g$), commutativity (\texttt{mul\_comm}), and the homomorphism law into one ``group object'' certificate. With $g^{-1} = g$ for all $g$, this is the abelian group $\mathbb{Z}/2 \times \mathbb{Z}/2$, proven by finite exhaustion over the $64$ triples.
\end{theorem}
\begin{proposition}[The internal $\mathbb{Z}/2$ gauge]\label{thm:gauge}\lean{Obsmetrologia.Positing.Gauge.K4.inGauge}\leanok\uses{thm:k4-isgroup}
\texttt{Positing.Gauge.K4.inGauge} marks the subgroup $\langle\sim\rangle = \{\texttt{id}, \texttt{neg}\} \cong \mathbb{Z}/2$ --- the internal gauge of FOUR. Its action on valuations is \texttt{Positing.Gauge.K4.gaugeAct}, where the nontrivial element \texttt{neg} acts as pointwise negation; \texttt{Positing.Gauge.K4.gauge\_closed} proves the subgroup is closed under multiplication.
\end{proposition}
\begin{theorem}[Gauge stability and movement]\label{thm:gauge-stable}\lean{Obsmetrologia.Positing.Gauge.K4.paraconsistent_gauge_stable}\leanok\uses{thm:gauge, thm:valuation}
\texttt{Positing.Gauge.K4.paraconsistent\_gauge\_stable} proves that if every cell of a valuation is $\bot_k$ or $\top_k$, the whole $\mathbb{Z}/2$ gauge fixes it: paraconsistent points are gauge-stable. Dually \texttt{Positing.Gauge.K4.classical\_gauge\_moves} proves a classical $t$-cell is moved by the gauge generator. The gauge fixes the inconsistency side and sign-flips the truth side --- a $\mathbb{Z}/2$ symmetry that distinguishes paraconsistent from classical content.
\end{theorem}

\subsection{The dihedral extension $D_4 = K_4 \sqcup K_4\cdot\sigma$}

\begin{proposition}[The order-swap involution $\sigma$]\label{thm:swap-kle-tle}\lean{Obsmetrologia.Charge.swap, Obsmetrologia.Charge.swap_kle_tle}\leanok\uses{thm:four-bilattice}
\texttt{Charge.swap}, written $\sigma$, is the fifth symmetry of FOUR that belongs to neither
order's automorphism group: in slot coordinates $\sigma\langle a,b\rangle = \langle a, \neg
b\rangle$, distinct from $\sim$ and $-$. \texttt{Charge.swap\_kle\_tle} proves $\sigma$ is an
order-isomorphism $(\mathrm{Charge},\le_k) \cong (\mathrm{Charge},\le_t)$: it carries the knowledge
order onto the truth order and back, exchanging the two orders rather than preserving or reversing
either one.
\end{proposition}
\begin{proposition}[The carrier $D_8$]\label{thm:d8-carrier}\lean{Obsmetrologia.D4.D8}\leanok\uses{thm:k4-group, thm:swap-kle-tle}
\texttt{D4.D8} is the eight-constructor inductive $\{\texttt{id}, \texttt{neg}, \texttt{conf},
\texttt{negConf}, \texttt{swap}, \texttt{negSwap}, \texttt{confSwap}, \texttt{negConfSwap}\}$: the
$K_4$ coset extended by the order-swap involution $\sigma$ (\texttt{Charge.swap}) and its three
$K_4$-composites. \texttt{D4.D8.act\_extends\_K4} proves the first four elements act exactly as
\texttt{Positing.Gauge.K4.act} does --- $D_8$ extends $K_4$, it does not redefine it.
\end{proposition}
\begin{theorem}[Faithful and a homomorphism]\label{thm:d8-faithful-hom}\lean{Obsmetrologia.D4.D8.faithful, Obsmetrologia.D4.D8.act_mul}\leanok\uses{thm:d8-carrier}
\texttt{D4.D8.faithful} proves the eight \texttt{act}-images are pairwise distinct (any two
elements disagree at some charge); \texttt{D4.D8.act\_mul} proves \texttt{act} is a group
homomorphism, so the $8\times 8$ Cayley table \texttt{D4.D8.mul} is exactly composition of the
underlying \texttt{Charge} maps.
\end{theorem}
\begin{theorem}[$D_8$ is the dihedral group of order 8]\label{thm:d4-is-dihedral}\lean{Obsmetrologia.D4.d4_is_dihedral}\leanok\uses{thm:d8-faithful-hom}
\texttt{D4.d4\_is\_dihedral} bundles: the group laws (associativity, two-sided identity, two-sided
inverses, with the order-4 rotations \texttt{negSwap}/\texttt{confSwap} mutually inverse and every
other element self-inverse); the dihedral relation $\sigma\circ{\sim}\circ\sigma = -$; the
presentation $\langle r, s \mid r^4 = s^2 = 1,\, srs = r^{-1}\rangle$ with $r := \texttt{negSwap}$,
$s := \texttt{swap}$; and non-abelian-ness ($\texttt{neg}\cdot\texttt{swap} \ne
\texttt{swap}\cdot\texttt{neg}$, the two distinct mutually-inverse order-4 elements pinning the
group to $D_4$ rather than $\mathbb{Z}/4\times\mathbb{Z}/2$ or $(\mathbb{Z}/2)^3$).
\end{theorem}
\begin{theorem}[$D_4 = \mathrm{Aut}(\mathrm{FOUR}; \{\le_k, \le_t\})$]\label{thm:d4-is-aut-four}\lean{Obsmetrologia.D4.d4_is_aut_four}\leanok\uses{thm:d4-is-dihedral}
\texttt{D4.d4\_is\_aut\_four} proves the semantic content: the $K_4$ coset preserves each order
setwise (never exchanging $\le_k$ and $\le_t$, via $K_4$'s existing order-effect lemmas), the
$\sigma$-coset exchanges $\le_k \leftrightarrow \le_t$ (via \texttt{Charge.swap\_kle\_tle}), and the
coset-membership bit is a homomorphism $D_8 \to \mathbb{Z}/2$ (\texttt{D4.D8.coset\_hom}), so $K_4$
is normal of index 2 --- the stabilizer of the \emph{ordered} pair $(\le_k, \le_t)$, with $D_4$ the
stabilizer of the \emph{unordered} pair $\{\le_k, \le_t\}$, the full dihedral symmetry of the
bilattice square.
\end{theorem}

\subsection{The finite/transfinite algebra/coalgebra charger}

\begin{proposition}[The charger is a new $K_4$ face]\label{thm:finite-infinite-charger-k4}\lean{Obsmetrologia.AperiodicOrder.Tiling.Charger.ofK4_toK4}\leanok\uses{thm:k4-group, thm:four-regime-dichotomy}
\texttt{AperiodicOrder.Tiling.Charger.Cell} is the product of two binary axes:
algebra/coalgebra and finite/transfinite.  \texttt{Charger.toK4} reads the four
cells as \(\texttt{id},\texttt{neg},\texttt{conf},\texttt{negConf}\), and
\texttt{Charger.ofK4\_toK4} with \texttt{Charger.toK4\_ofK4} proves the
round trips.  This is not a new primitive four-point object; it is another
presentation of the existing Klein-four carrier.
\end{proposition}

\begin{theorem}[The two charger axes commute]\label{thm:finite-infinite-charger-axes}\lean{Obsmetrologia.AperiodicOrder.Tiling.Charger.axes_commute}\leanok\uses{thm:finite-infinite-charger-k4}
\texttt{Charger.arrowDual} toggles algebra and coalgebra, while
\texttt{Charger.finitenessDual} toggles finite and transfinite.  The theorem
\texttt{Charger.axes\_commute} proves these two \(\mathbb{Z}/2\) actions commute,
and \texttt{Charger.arrowDual\_toK4}, \texttt{Charger.finitenessDual\_toK4},
and \texttt{Charger.composite\_toK4} identify the toggles with left
multiplication by the three nontrivial elements of \(K_4\).  The existing
\texttt{Charge.swap\_swap} order-swap duality is recorded only as an
arrow-reversal-shaped witness; it is distinct from the original charge \(K_4\)
elements.
\end{theorem}

\begin{theorem}[The off-diagonal bridges]\label{thm:finite-infinite-charger-bridges}\lean{Obsmetrologia.AperiodicOrder.Tiling.Charger.bridge_cells}\leanok\uses{thm:finite-infinite-charger-axes, thm:weave-draft-sofic, thm:jeandel-rao-substitution-tower}
\texttt{Charger.bridge\_cells} identifies the two off-diagonal mechanisms:
\texttt{Bridge.soficCover} lands in the finite-coalgebra cell and
\texttt{Bridge.directedColimit} lands in the transfinite-algebra cell.
The concrete witnesses are \texttt{Charger.fib\_deep\_sofic\_lands\_finite\_coalgebra},
which charges the decorated Fibonacci sofic cover, and
\texttt{Charger.colimit\_cocone\_lands\_transfinite\_algebra}, which keeps the
directed-colimit cocone law while charging it as transfinite algebra.  The
charger therefore provides a routing use: analytic/coinductive difficulty can be
approached through the finite-coalgebra sofic cover or the transfinite-algebra
colimit, rather than by confusing the two diagonals.
\end{theorem}
